\section{Immutable Parent Pointers}
\label{sec:immutable-parent-pointers}

The immutable parent pointer is the foundational structural primitive of the Recursive Spawning Protocol. It is the single mechanism that renders the delegation chain a runtime-verifiable artifact rather than a forensic reconstruction, makes chain integrity a structural guarantee rather than a policy assertion, and makes all modes of autonomous drift architecturally impossible regardless of the operational urgency, architectural complexity, or adversarial sophistication of any machine actor in the chain. This section specifies the parent pointer formally, defines the chain traversal operator using the recursive formulation, establishes the base case for the root human anchor, specifies the immutability enforcement mechanism through the Fact Layer write protocol, describes the Purpose Layer's spawn authorization role, and provides the complete chain traversal algorithm with correctness arguments.

% ------------------------------------------------------------------------------%
\subsection{Formal Definition: ParentPointer}
\label{sec:parent-pointer-definition}

\begin{definition}[Parent Pointer]
\label{def:parent-pointer}
Let $\mathcal{N}$ be the set of all agent nodes in a \bxthree{} system implementing the Recursive Spawning Protocol. For any node $c \in \mathcal{N}$ spawned via a valid RSP spawn event, the relation $\ParentPointer{p}{c}$ is established exactly once at node provisioning via a Fact Layer atomic write, and satisfies four governing properties:

\begin{enumerate}
    \item[(P1) \textbf{Uniqueness}] For every child node $c \in \mathcal{N}$, there exists exactly one node $p \in \mathcal{N}$ such that $\ParentPointer{p}{c}$ holds at all times after provisioning. No child may have multiple simultaneous parent pointers. No mechanism exists by which an established pointer may be overwritten or redirected.

    \item[(P2) \textbf{Immutability}] After the provisioning write confirming $\ParentPointer{p}{c}$ is recorded in the Fact Layer ledger, no operation exists --- from any source, with any privilege level, under any system condition --- that can modify or delete $\ParentPointer{p}{c}$. The pointer is permanent for the operational lifetime of $c$. This is not a policy enforced by access control; it is a structural property of the protocol: the modification operation is not defined.

    \item[(P3) \textbf{Directionality}] The parent pointer is an intrinsic attribute of the child node, not a reference held by the parent. The child $c$ carries $\ParentPointer{p}{c}$ as an immutable attribute in its own Fact Layer state record. The parent $p$ holds no mutable reference to $c$. A parent that has terminated, crashed, or been decommissioned does not affect $c$'s ability to reference its parent. The chain is traversable from any node outward to its complete lineage independent of the operational state of any other node.

    \item[(P4) \textbf{Human Anchor Termination}] The root of any RSP chain is a human principal node $h$ for which $\ParentPointer{\bot}{h}$ holds. The null parent marker $\bot$ indicates that $h$ is the ultimate delegation origin. No machine actor in the RSP model may serve as a chain root.
\end{enumerate}
\end{definition}

The distinction between P2-immutability and access-control-based immutability is architectural. An access-controlled pointer can be modified by any principal with sufficient privileges; immutability is a policy enforceable until credential theft, privilege escalation, or administrative compromise circumvents it. The Fact Layer write protocol makes immutability a structural property: there is no modification pathway defined, so no amount of privilege elevation can produce a valid write to an established $\ParentPointer{p}{c}$. The pointer is not protected --- it is structurally incapable of being changed.

\medskip
\noindent\textbf{Notation.} We write $\ParentPointer{p}{c}$ to denote the parent pointer relation between parent $p$ and child $c$. We write $\alpha(c) = p$ and $\parent(c) = p$ interchangeably as functional notation. The expression $\alpha(c) = \bot$ (or equivalently $\parent(c) = \bot$) denotes that node $c$ has no parent --- it is the root human anchor. The functional notation $\alpha^{(k)}(a)$ denotes the $k$-fold application of $\alpha$, and $\alpha^{(0)}(a) = a$.

\begin{dfn}[Spawn Event]
\label{dfn:spawn-event}
A \emph{spawn event} is the 5-tuple
\begin{equation}
\mathrm{spawn}(c,\ p,\ \sigma_c,\ t,\ \phi)
\end{equation}
where $c$ is the newly created child node, $p \in \mathcal{N}$ is the authorizing parent node, $\sigma_c$ is the authorized execution scope, $t$ is the wall-clock timestamp, and $\phi$ is the cryptographic signature produced by the Purpose Layer over the spawn request. The Fact Layer atomically: (1) creates node $c$, (2) sets $\ParentPointer{p}{c}$, (3) sets the node's authorized scope to $\sigma_c$, (4) seals the memory envelope, and (5) appends the event to the forensic ledger $\mathcal{L}$.
\end{dfn}

The atomic coupling of authorization (Purpose Layer) with assignment (Fact Layer) prevents a parent from claiming to have spawned a node whose parent pointer was set to a different actor, and prevents a child from falsifying its own ancestry after provisioning.

\begin{prop}[Immutability of $\ParentPointer{p}{c}$]
\label{prop:parent-pointer-immutable}
For any child node $c$ with parent $p$, once the spawn event $\mathrm{spawn}(c, p, \sigma_c, t, \phi)$ has been recorded in the forensic ledger $\mathcal{L}$, the relation $\ParentPointer{p}{c}$ is fixed for the operational lifetime of $c$. No subsequent operation --- including operations issued by $p$, by any Purpose Layer actor, by the Bounds Engine, or by $c$ itself --- can alter $\ParentPointer{p}{c}$.
\end{prop}

\begin{pf}
The Fact Layer write protocol defines no code path for processing a modification request to an established $\ParentPointer{p}{c}$ (W1). All modification requests are rejected at the protocol level before the Fact Layer attempts to execute them (W2). Since there is no operation definition for modification, no actor --- regardless of privilege level --- can issue a valid modification. The atomic provisioning property (W4) ensures $\ParentPointer{p}{c}$ and its warrant $\phi$ are created simultaneously and sealed. Therefore $\ParentPointer{p}{c}$ cannot change after provisioning. ∎
\end{pf}

% ------------------------------------------------------------------------------%
\subsection{The Chain Operator}
\label{sec:chain-operator}

The chain operator $\Chain{\cdot}$ constructs the complete delegation lineage from any node to the human anchor by recursively applying the parent pointer function. This operator is the primary tool for chain verification, audit, and accountability attribution.

\begin{dfn}[Chain Operator]
\label{def:chain-operator}
For any agent node $a \in \mathcal{N}$ in a Recursive Spawning chain, the \emph{chain} $\Chain{a}$ is defined recursively as:
\begin{equation}
\Chain{a} =
\begin{cases}
\{\,a,\ \ParentPointer{\parent(a)}{a}\,\} \cup \Chain{\parent(a)} & \text{if } \alpha(a) \neq \bot \\[10pt]
\{\,a,\ h(a)\,\} & \text{if } \alpha(a) = \bot \quad \text{(root human anchor)}
\end{cases}
\label{eq:chain-recursive}
\end{equation}
Equivalently, expanding the recursion to reveal the full lineage:
\begin{equation}
\Chain{a} = \bigl\{\, a,\ \alpha(a),\ \alpha^{(2)}(a),\ \alpha^{(3)}(a),\ \ldots,\ \alpha^{(k)}(a) \,\bigr\}
\label{eq:chain-expanded}
\end{equation}
where $k$ is the depth of node $a$ (the number of successive parent applications required to reach the root), $\alpha^{(j)}(a)$ denotes the $j$-fold application of $\alpha$, and $\alpha^{(k)}(a)$ is the unique root node $n_0$ satisfying $\alpha(n_0) = \bot$ and $h(n_0)$.
\end{dfn}

The base case of the recursion is the root human anchor, where $\alpha(n_0) = \bot$ and the node carries the human anchor designation $h(\cdot)$. This is the only valid chain root in RSP: no machine actor may serve as a chain origin, and no node may have $\alpha(n) = \bot$ unless it carries the $h(\cdot)$ designation. The non-collapsing human anchor property ensures $h(n) = h(n_0)$ for all nodes in all chains spawned from $n_0$. The chain $\Chain{a}$ is always finite: by the RSP termination property, every chain has depth at most $D_{\max}$, so the recursion terminates within at most $D_{\max} + 1$ successive applications of $\alpha$.

The chain operator satisfies three critical properties that make it suitable for runtime verification:

\begin{itemize}
    \item[(C1) \textbf{Determinism}] For any node $a$, $\Chain{a}$ contains exactly one element for each ancestor on the delegation path. There is no alternative chain, no ambiguous parentage, and no non-deterministic branching. The parent pointer function $\alpha$ is a total function with unique values at every node.

    \item[(C2) \textbf{Stability}] The immutability of $\alpha$ guarantees that $\Chain{a}$ at any runtime point $t$ is identical to $\Chain{a}$ at any prior runtime point $t' < t$. The chain does not change as a function of time or operational state. This is what makes the chain a runtime artifact rather than a forensic reconstruction.

    \item[(C3) \textbf{Verifiability}] Any observer with access to the Fact Layer ledger can independently compute $\Chain{a}$ and obtain the same result. Chain integrity does not depend on trusting any node in the chain to report its own ancestry accurately.
\end{itemize}

\begin{corol}[Chain Terminates at Human Anchor]
\label{corol:chain-terminates-human}
For any node $a \in \mathcal{N}$, $\Chain{a}$ contains a unique human principal node $h(n_0)$ and no machine actor serves as a chain terminus.
\end{corol}

\begin{pf}
By definition, $\Chain{a}$ is constructed by repeated application of $\alpha$ until $\alpha(n) = \bot$. By property (P4), only a node carrying the $h(\cdot)$ designation may have $\alpha(n) = \bot$. No machine actor carries the $h(\cdot)$ designation. Therefore $\Chain{a}$ terminates at a human principal and at no other node type. ∎
\end{pf}

% ------------------------------------------------------------------------------%
\subsection{Immutability Enforcement: Fact Layer Write Protocol}
\label{sec:fact-layer-write-protocol}

The immutability of $\ParentPointer{p}{c}$ is not a policy choice or an access-control configuration. It is a structural property of the Recursive Spawning Protocol enforced by the Fact Layer write protocol, which defines what operations are valid and what operations are rejected at the protocol level, before any execution context is entered.

The Fact Layer maintains four write properties that collectively make $\ParentPointer{p}{c}$ immutable after provisioning:

\begin{enumerate}
    \item[(W1) \textbf{No modification operation defined.}] The Fact Layer write protocol contains no code path for processing a modification request to an established $\ParentPointer{p}{c}$ relation. The operation does not exist in the instruction set. This is distinct from a permission denied response: the protocol treats modification requests as malformed inputs, not as authorized operations that happen to fail on access control checks.

    \item[(W2) \textbf{Protocol-level rejection of modification.}] All modification requests --- regardless of source, privilege level, or cryptographic credential --- are rejected at the protocol level before the Fact Layer attempts to execute them. There is no privilege level, administrative role, or emergency override that can produce a valid modification operation.

    \item[(W3) \textbf{Sealed memory envelope.}] The parent pointer field $\ParentPointer{p}{c}$ is encrypted at rest using a Fact Layer-only symmetric key. Even a compromised system component cannot read $\ParentPointer{p}{c}$ because it never has access to the decryption key. The envelope carries an HMAC-SHA-256 signature computed over its contents; any modification to the ciphertext invalidates the signature. The Fact Layer decrypts the envelope only when reading $\ParentPointer{p}{c}$ to service a chain-traversal or audit request, and re-encrypts and re-signs after each authorized read, preventing replay and corruption attacks.

    \item[(W4) \textbf{Atomic provisioning.}] The parent pointer $\ParentPointer{p}{c}$ and its Purpose Layer warrant $\phi$ are created as a single atomic Fact Layer write. There is no temporal window in which the pointer exists without authorization, or in which authorization exists without a corresponding pointer. This eliminates the class of exploits where a parent pointer is first established and then retroactively authorized.
\end{enumerate}

\medskip
\noindent\textbf{Why access control is insufficient.} The distinction between access-control-based immutability and protocol-level immutability is the distinction between a promise and a guarantee. In an access-control model, immutability is enforced by restricting which principals may issue modification operations. If an attacker acquires sufficient privileges --- through credential theft, insider compromise, or software vulnerability --- the access control barrier is circumvented and immutability is violated. The immutability was a policy, enforceable until it was not.

In the Fact Layer write protocol, the immutability of $\ParentPointer{p}{c}$ is a property of the protocol itself. There is no operation defined for modifying $\ParentPointer{p}{c}$ after provisioning, so no amount of privilege elevation, credential compromise, or software vulnerability can produce a valid modification operation. The protocol does not enforce immutability --- it makes immutability the only valid state transition. The distinction is architectural, not a matter of degree.

This is the precise failure mode the Fact Layer write protocol eliminates: retroactive manipulation of chain topology for the purpose of accountability laundering. Without protocol-level immutability, an adversarial actor could re-parent a node to a favorable parent after the fact, obscuring the original authorization chain. With protocol-level immutability, this manipulation is structurally impossible.

\medskip
\noindent\textbf{Four enforcement properties.} The immutability guarantee exhibits four properties that distinguish it from access-control-based approaches:

\begin{enumerate}
    \item[(E1) \textbf{No override path.}] There is no configuration flag, emergency pathway, or administrative override capable of modifying $\ParentPointer{p}{c}$ after provisioning. Access control systems have override paths; the Fact Layer write protocol does not.

    \item[(E2) \textbf{Failure-state resilience.}] The immutability holds even under system conditions that disable other enforcement mechanisms --- including catastrophic failure, kernel compromise, or memory corruption in other system components. The parent pointer is stored in non-volatile, sealed Fact Layer memory. A node that experiences power loss or hardware damage resumes with its parent pointer intact.

    \item[(E3) \textbf{Source-equivalence of enforcement.}] All modification requests --- from any source --- receive identical treatment: protocol-level rejection. The Fact Layer does not distinguish between principals. A request from the Purpose Layer, the Bounds Engine, an authenticated administrator, or an external adversarial actor receives the same response.

    \item[(E4) \textbf{Atomic warrant-pointer creation.}] The parent pointer and its Purpose Layer warrant are created as a single atomic Fact Layer write. There is no temporal window in which $\ParentPointer{p}{c}$ is established without a matching warrant, or in which a warrant exists without a corresponding parent pointer.
\end{enumerate}

% ------------------------------------------------------------------------------%
\subsection{Spawn Authorization: Purpose Layer Role}
\label{sec:purpose-layer-authorization}

While existing $\ParentPointer{p}{c}$ relations are immutable, the system must be able to grow the accountability chain through new spawn events. The creation of a new $\ParentPointer{p}{c}$ is the sole mechanism by which the chain extends, and it occurs exactly once per spawn event --- at node provisioning, atomically and with mandatory Purpose Layer authorization.

A spawn event is valid under RSP only when three conditions are satisfied simultaneously:

\begin{enumerate}
    \item[(S1) \textbf{Human authorization.}] The parent node $p$ requesting the spawn holds a Purpose Layer warrant $\phi$ issued by a human principal $h(p)$ confirming that $h(p)$ has authorized the creation of the child node $c$ with authorized scope $\sigma_c$. The warrant $\phi$ is a cryptographic signature over the spawn parameters $(c, p, \sigma_c, t)$ produced using the human principal's private key and verifiably tied to $h(p)$'s public key in the Fact Layer registry.

    \item[(S2) \textbf{Scope refinement.}] The authorized scope $\sigma_c$ of the child node $c$ is a descendant refinement of the parent node's authorized scope $R(p)$. Formally: $\sigma_c \subseteq R(p)$. The child may not operate in any domain outside the parent's authorized scope. Lateral expansion of operational authority beyond the parent chain's authorized intent is structurally prohibited.

    \item[(S3) \textbf{Depth availability.}] The depth of the proposed child node $c$, defined as $\mathrm{depth}(c) = \mathrm{depth}(p) + 1$, does not exceed the RSP depth bound $D_{\max}$. The parent node $p$ is responsible for checking this condition before issuing the spawn request. The Bounds Engine independently verifies depth compliance as a precondition for Fact Layer provisioning.
\end{enumerate}

The Purpose Layer's role in spawn authorization is exclusive: it carries the intent and judgment that the Fact Layer enforces. The Fact Layer cannot issue a parent pointer to a node that lacks a Purpose Layer warrant; the Bounds Engine cannot extend the depth bound to accommodate a spawn that would exceed $D_{\max}$. Each layer performs its own gate check, and all three gates must pass before a parent pointer is established. This mandatory three-layer concert is what distinguishes RSP authorization from policy-based authorization schemes that rely on a single enforcement point.

% ------------------------------------------------------------------------------%
\subsection{Chain Traversal Algorithm}
\label{sec:chain-traversal-algorithm}

This section presents the runtime procedure for constructing the delegation chain from any node to its human anchor. It is provided here in full algorithmic detail with correctness arguments for termination, correctness, uniqueness, and immutability guarantees.

\medskip
\noindent\textbf{Algorithm: Chain-Traverse.}

\begin{algorithm}
\caption{Chain-Traverse($n$) --- Build the delegation chain from node $n$ to the human anchor}
\label{alg:chain-traverse}
\begin{algorithmic}[1]
\REQUIRE $n$ is a provisioned RSP node with $\alpha$ defined
\ENSURE $\Chain{n}$ --- the ordered set of all ancestor nodes including the human anchor

\STATE $chain \leftarrow \langle\ \rangle$ \COMMENT{Initialize empty ordered list}
\STATE $current \leftarrow n$ \COMMENT{Start at the query node}

\WHILE{$\mathrm{true}$}
    \STATE $\mathrm{append}(chain,\ current)$ \COMMENT{Add current node to chain}
    \IF{$\alpha(current) = \bot$}
        \STATE $\mathrm{append}(chain,\ h(current))$ \COMMENT{Add human anchor and terminate}
        \STATE \textbf{break}
    \ENDIF
    \STATE $current \leftarrow \alpha(current)$ \COMMENT{Move to parent via $\alpha$}
\ENDWHILE

\RETURN $chain$
\end{algorithmic}
\end{algorithm}

\medskip
\noindent\textbf{Correctness arguments.}

\textit{Termination.} The algorithm terminates because RSP chains are depth-bounded by construction. Each iteration of the while-loop moves $current$ to the parent node $\alpha(current)$, which is strictly closer to the root along the chain. Since the chain has finite depth at most $D_{\max} + 1$, the loop executes at most $D_{\max} + 1$ iterations before $\alpha(current) = \bot$ is reached. This argument holds regardless of the current operational state of any node in the chain, including nodes that may have terminated or entered error states, because the traversal only reads immutable $\alpha$ values.

\textit{Correctness of the returned set.} By Definition~\ref{def:chain-operator}, $\Chain{a}$ is the set containing $a$, $\alpha(a)$, $\alpha(\alpha(a))$, and so on, until the root. The algorithm produces exactly this sequence: it starts at $a$, appends each node, moves to its parent via $\alpha$, and stops when $\alpha(current) = \bot$ is reached. The output is identical to $\Chain{a}$ by construction.

\textit{Uniqueness of result.} The algorithm produces a unique result for any input node because $\alpha$ is a total function (defined for every provisioned node) and is injective on the chain direction (every node has exactly one parent, and the root has no parent). There is no branching, no alternative path, and no non-deterministic choice at any step.

\textit{Immutability guarantee.} The algorithm reads $\alpha$ values but never writes them. The traversal is read-only and does not interact with the Fact Layer write protocol. The algorithm can be executed at any runtime point without affecting the immutability of the chain it traverses.

\medskip
\noindent\textbf{Algorithm: Chain-Verify.}

\begin{algorithm}
\caption{Chain-Verify($n$) --- Verify the authorization chain from node $n$ to human anchor}
\label{alg:chain-verify}
\begin{algorithmic}[1]
\REQUIRE $n$ is a provisioned RSP node
\ENSURE $\mathrm{Boolean}$ indicating whether the chain is a structurally valid RSP chain

\STATE $chain \leftarrow \mathrm{Chain-Traverse}(n)$
\STATE $m \leftarrow \mathrm{len}(chain)$

\STATE \COMMENT{\textit{V1: Verify chain terminates at human anchor}}
\IF{$\alpha(chain[m-1]) \neq \bot$}
    \STATE \RETURN $\mathrm{false}$ \COMMENT{Root is not $\bot$ --- malformed chain}
\ENDIF
\IF{$\neg h(chain[m-1])$}
    \STATE \RETURN $\mathrm{false}$ \COMMENT{Root is not a designated human principal}
\ENDIF

\STATE \COMMENT{\textit{V2: Verify each spawn event has a Purpose Layer warrant}}
\FOR{$i \leftarrow 0$ \TO $m-2$}
    \STATE $child \leftarrow chain[i]$
    \STATE $parent \leftarrow chain[i+1]$
    \IF{$\neg \mathrm{warrant\_exists}(child,\ parent)$}
        \STATE \RETURN $\mathrm{false}$ \COMMENT{Missing Purpose Layer warrant for spawn}
    \ENDIF
\ENDFOR

\STATE \COMMENT{\textit{V3: Verify scope refinement at each link}}
\FOR{$i \leftarrow 0$ \TO $m-2$}
    \STATE $child \leftarrow chain[i]$
    \STATE $parent \leftarrow chain[i+1]$
    \IF{$\neg (R(child) \subset R(parent))$}
        \STATE \RETURN $\mathrm{false}$ \COMMENT{Scope refinement constraint violated}
    \ENDIF
\ENDFOR

\STATE \COMMENT{\textit{V4: Verify depth bound compliance at each node}}
\FOR{$i \leftarrow 0$ \TO $m-1$}
    \IF{$\mathrm{depth}(chain[i]) > D_{\max}$}
        \STATE \RETURN $\mathrm{false}$ \COMMENT{Depth bound exceeded at node}
    \ENDIF
\ENDFOR

\RETURN $\mathrm{true}$
\end{algorithmic}
\end{algorithm}

Chain-Verify returns \texttt{true} if and only if all four verification conditions hold simultaneously. V1 confirms the chain terminates at a human anchor (not at a machine actor or orphaned node). V2 confirms every spawn was authorized by the Purpose Layer (not self-bootstrapped by a machine actor). V3 confirms scope refinement was respected at every link (not lateral expansion beyond authorized intent). V4 confirms the depth bound was respected throughout (not approaching infinity). A chain that passes Chain-Verify is a fully authorized delegation chain terminating at a human principal.

The algorithm is deterministic and produces the same result regardless of which node in the chain initiates verification. There is no self-serving interpretation of authorization that a drifted or orphaned node could produce, because Chain-Verify inspects the complete chain topology and warrants in the Fact Layer ledger, not the node's own claims about its authorization.

\begin{prop}[Chain-Verify Soundness]
\label{prop:chain-verify-soundness}
If Chain-Verify($n$) returns \texttt{true}, then $\Chain{n}$ is a valid RSP delegation chain terminating at a human principal, and all spawn events in the chain satisfy the Purpose Layer authorization conditions (S1), (S2), and (S3).
\end{prop}

\begin{prop}[Chain-Verify Completeness]
\label{prop:chain-verify-completeness}
If $\Chain{n}$ is a valid RSP delegation chain terminating at a human principal, then Chain-Verify($n$) returns \texttt{true}.
\end{prop}

\medskip
\noindent\textbf{Computational cost.} Chain-Traverse runs in $O(d)$ time where $d$ is the chain depth (at most $D_{\max} + 1$). Chain-Verify runs in $O(m)$ time for the traversal plus $O(m)$ for each of the three verification loops, yielding $O(m)$ total where $m$ is chain length (also bounded by $D_{\max} + 1$). The algorithms are linear in chain length and independent of the total number of nodes in the system.

The depth bound $D_{\max}$ provides a constant-time upper bound on chain traversal cost: the worst-case number of pointer dereferences required to reach the human anchor from any node is $D_{\max} + 1$, independent of the total number of nodes in the system. Mutable parent pointers would allow the chain to be extended retroactively, making the effective depth potentially unbounded and the worst-case escalation time unbounded as well. Immutability is what makes the depth bound a structural guarantee, not merely a policy.

% ------------------------------------------------------------------------------%
\subsection{Summary}
\label{sec:immutable-parent-pointers-summary}

The immutable parent pointer is the load-bearing constraint of the Recursive Spawning Protocol. Without it, the chain is mutable --- subject to retroactive modification by actors with sufficient privilege --- and the RSP's correctness properties (drift prevention, chain integrity, termination) collapse to policy conventions rather than structural guarantees.

The immutability is architectural --- not a stronger access control policy --- because it is enforced by the Fact Layer write protocol, which defines no valid modification operation for $\ParentPointer{p}{c}$ after provisioning. This is distinct from access-control-based immutability in four fundamental ways: there is no override path (E1); immutability holds under all failure states (E2); all sources receive equivalent enforcement treatment regardless of privilege level (E3); and the parent pointer and its Purpose Layer warrant are atomically created (E4).

The chain operator $\Chain{a}$, defined recursively as $\ParentPointer{\parent(a)}{a} \cup \Chain{\parent(a)}$ with base case at the root human anchor where $\alpha(n) = \bot$, is uniquely determined by the immutable parent pointers in the chain (Corollary~\ref{corol:chain-terminates-human}). The chain traversal algorithm (Algorithm~\ref{alg:chain-traverse}) computes this operator correctly and terminates. Chain-Verify (Algorithm~\ref{alg:chain-verify}) provides the runtime verification procedure for chain structural validity (Propositions~\ref{prop:chain-verify-soundness} and~\ref{prop:chain-verify-completeness}). Together, these mechanisms make the delegation chain an immutable, verifiable, and auditable runtime artifact --- the foundation on which all other RSP guarantees rest.
